% Documento principal. No escribas contenido aquí.
%   - Lo estructural (título, problemas, objetivos, hipótesis, variables, matrices,
%     cronograma, presupuesto, carátula) se edita en tesis.yaml.
%   - El orden de capítulos y la numeración oficial UNTELS los arma
%     generado/estructura.tex según meta.modo (plan | tesis) y meta.con_hipotesis.
%   - El texto de cada sección vive en capitulos/, preliminares/ y anexos/.
% Compilar: scripts/compilar.sh (o .ps1) desde la skill `tesis`, que primero corre
% generar.py. En Overleaf compila tal cual, porque generado/ está versionado.
\documentclass{untels}

% GENERADO por generar.py desde tesis.yaml. No editar a mano.
\modoplan
\conhipotesistrue
\ConfigurarInterlineado{1.5}
\ConfigurarAlineacion{justificado}
\newcommand{\TituloTesis}{SISTEMA WEB CON CHATBOT DE WHATSAPP E INTELIGENCIA ARTIFICIAL PARA LA GESTIÓN DE ATENCIÓN AL CLIENTE EN UNA EMPRESA DE SERVICIOS}
\newcommand{\AutorPrincipal}{Nombre Nombre Apellido Apellido}
\newcommand{\AutorPrincipalDNI}{00000000}
\newcommand{\Asesor}{Mg. Apellido Apellido, Nombre}
\newcommand{\AsesorORCID}{0000-0000-0000-0000}
\newcommand{\Universidad}{UNIVERSIDAD NACIONAL TECNOLÓGICA DE LIMA SUR}
\newcommand{\Facultad}{FACULTAD DE INGENIERÍA Y GESTIÓN}
\newcommand{\Escuela}{ESCUELA PROFESIONAL DE INGENIERÍA DE SISTEMAS}
\newcommand{\Carrera}{INGENIERÍA DE SISTEMAS}
\newcommand{\TituloProfesional}{INGENIERO DE SISTEMAS}
\newcommand{\Ciudad}{Villa El Salvador}
\newcommand{\Anio}{2026}
\newcommand{\PalabrasClave}{chatbot, WhatsApp, inteligencia artificial, atención al cliente}
\newcommand{\Keywords}{chatbot, WhatsApp, artificial intelligence, customer service}
\newcommand{\DelimEspacial}{EJEMPLO: la investigación se desarrollará en el área de atención al cliente de una empresa de servicios ubicada en Villa El Salvador, Lima.}
\newcommand{\DelimTemporal}{EJEMPLO: la investigación se desarrollará entre marzo y diciembre de 2026; la pre-prueba corresponde a abril y la post-prueba a octubre.}
\hypersetup{pdftitle={SISTEMA WEB CON CHATBOT DE WHATSAPP E INTELIGENCIA ARTIFICIAL PARA LA GESTIÓN DE ATENCIÓN AL CLIENTE EN UNA EMPRESA DE SERVICIOS},pdfauthor={Nombre Nombre Apellido Apellido}}
          % \modoplan|\modotesis, \conhipotesistrue|false, datos de carátula
\addbibresource{bib/referencias.bib}

\begin{document}
% GENERADO por generar.py desde tesis.yaml. No editar a mano.
% Estructura oficial UNTELS: plan con hipótesis.
\pagenumbering{arabic}
% GENERADO por generar.py desde tesis.yaml. No editar a mano.
\thispagestyle{empty}
\begin{center}
\fontsize{14}{17}\selectfont\setstretch{1}
{\bfseries UNIVERSIDAD NACIONAL TECNOLÓGICA DE LIMA SUR}\\[2pt]
{\bfseries FACULTAD DE INGENIERÍA Y GESTIÓN}\\[2pt]
{\bfseries CARRERA PROFESIONAL}\\[2pt]
INGENIERÍA DE SISTEMAS
\vfill
\includegraphics[width=4.6cm]{logo-untels}
\vfill
{\bfseries PLAN DE TESIS}
\vfill
\parbox{0.92\textwidth}{\centering ``SISTEMA WEB CON CHATBOT DE WHATSAPP E INTELIGENCIA ARTIFICIAL PARA LA GESTIÓN DE ATENCIÓN AL CLIENTE EN UNA EMPRESA DE SERVICIOS''}
\vfill
{\bfseries PARA OPTAR EL TÍTULO PROFESIONAL DE}\\[2pt]
INGENIERO DE SISTEMAS
\vfill
{\bfseries PRESENTADO POR:}\\[8pt]
\begin{tabular}{@{}ll@{}}
BACHILLER: & APELLIDO APELLIDO, NOMBRE NOMBRE \\
ASESOR: & Mg. APELLIDO APELLIDO, NOMBRE \\
\end{tabular}
\vfill
Villa El Salvador, 2026\\[2pt]
PERÚ
\end{center}
\clearpage

\tableofcontents
\listoftables
\listoffigures
\chapter{PLANTEAMIENTO DEL PROBLEMA}
% I. PLANTEAMIENTO DEL PROBLEMA. Lo escribe la skill tesis-problema.
% El título del capítulo lo pone generado/estructura.tex; aquí van las secciones.

\section{Descripción del problema}
% Realidad problemática: internacional -> nacional -> local/organización.
% Síntomas -> causas -> pronóstico -> control al pronóstico. Toda cifra o afirmación con cita.
% Sin subtítulos, en tercera persona; incluye motivación y estado del arte del problema.

\section{Formulación del problema}
% GENERADO por generar.py desde tesis.yaml. No editar a mano.
\subsection{Problema general}
¿En qué medida un sistema web con chatbot de WhatsApp e inteligencia artificial mejora la gestión de atención al cliente en una empresa de servicios?

\subsection{Problemas específicos}
\begin{enumerate}[label=\textbf{PE\arabic*:},leftmargin=1.27cm,itemsep=4pt]
\item ¿En qué medida un sistema web con chatbot de WhatsApp e inteligencia artificial reduce el tiempo promedio de primera respuesta al cliente?
\item ¿En qué medida un sistema web con chatbot de WhatsApp e inteligencia artificial incrementa la tasa de consultas resueltas sin intervención humana?
\end{enumerate}


\section{Objetivos de la investigación}
% GENERADO por generar.py desde tesis.yaml. No editar a mano.
\subsection{Objetivo general}
Determinar en qué medida un sistema web con chatbot de WhatsApp e inteligencia artificial mejora la gestión de atención al cliente en una empresa de servicios.

\subsection{Objetivos específicos}
\begin{enumerate}[label=\textbf{OE\arabic*:},leftmargin=1.27cm,itemsep=4pt]
\item Determinar en qué medida un sistema web con chatbot de WhatsApp e inteligencia artificial reduce el tiempo promedio de primera respuesta al cliente.
\item Determinar en qué medida un sistema web con chatbot de WhatsApp e inteligencia artificial incrementa la tasa de consultas resueltas sin intervención humana.
\end{enumerate}


\section{Delimitación de la investigación}
\subsection{Delimitación espacial}
\DelimEspacial

\subsection{Delimitación temporal}
\DelimTemporal

\section{Justificación del problema}
% Un \subsection por tipo, según corresponda: teórica, tecnológica, social, económica.
% Con cifras citadas.

\chapter{MARCO TEÓRICO}
% II. MARCO TEÓRICO. Lo escribe la skill tesis-marco.

\section{Antecedentes de la investigación}
\subsection{Antecedentes internacionales}
% Mínimo 5, de los últimos meta.antecedentes_ventana_anios años. Un párrafo por antecedente,
% \textcite{clave}, parafraseado y sin página: autor y año -> lugar -> objetivo -> diseño,
% población/muestra -> instrumentos -> resultados con cifras -> conclusión vinculada.

\subsection{Antecedentes nacionales}
% Mínimo 3; tesis o artículos peruanos recientes. Misma plantilla.

\section{Bases teóricas}
% Un \subsection por variable (VI y VD) con sus dimensiones e indicadores y fórmulas.
% Mínimo 3 definiciones por variable de autores distintos, integradas.
%
% Figura APA (el \caption va ANTES de la imagen):
% \begin{figure}[H]
%   \caption{Arquitectura del sistema propuesto}
%   \label{fig:arquitectura}
%   \centering\includegraphics[width=0.85\textwidth]{arquitectura}
%   \nota{Elaboración propia.}
% \end{figure}
%
% Tabla APA (solo líneas horizontales):
% \begin{table}[H]
%   \caption{Comparación de modelos de lenguaje}
%   \label{tab:modelos}
%   \begin{tabularx}{\textwidth}{@{}lYr@{}}
%     \toprule
%     \textbf{Modelo} & \textbf{Descripción} & \textbf{Costo} \\
%     \midrule
%     A & ... & 0.00 \\
%     \bottomrule
%   \end{tabularx}
%   \nota{Adaptado de \textcite{clave}.}
% \end{table}

\section{Definición de términos básicos}
% 10 a 20 términos en orden alfabético, cada uno citado:
% \textbf{Término.} Definición parafraseada \parencite{clave}.

\chapter{VARIABLES E HIPÓTESIS}
% III. VARIABLES E HIPÓTESIS (solo con hipótesis). Lo escribe la skill tesis-variables.

\section{Definición operacional de las variables}
% Uno o dos párrafos que presenten la VI y la VD y remitan a la tabla de operacionalización.
% GENERADO por generar.py desde tesis.yaml. No editar a mano.
\begin{landscape}
{\small\setlength{\tabcolsep}{4pt}
\begin{longtable}{L{2.4cm}L{3.4cm}L{3.4cm}L{2.4cm}L{2.8cm}L{3.3cm}L{1.5cm}L{2.0cm}}
\caption{Matriz de operacionalización de las variables}\label{tab:operacionalizacion}\\
\toprule
\textbf{Variable} & \textbf{Definición conceptual} & \textbf{Definición operacional} & \textbf{Dimensión} & \textbf{Indicador} & \textbf{Fórmula} & \textbf{Escala} & \textbf{Instrumento} \\
\midrule
\endfirsthead
\multicolumn{8}{@{}l}{\itshape (continuación)}\\
\toprule
\textbf{Variable} & \textbf{Definición conceptual} & \textbf{Definición operacional} & \textbf{Dimensión} & \textbf{Indicador} & \textbf{Fórmula} & \textbf{Escala} & \textbf{Instrumento} \\
\midrule
\endhead
\bottomrule
\endlastfoot
\textbf{VI:} Sistema web con chatbot de WhatsApp e inteligencia artificial & EJEMPLO: aplicación web que integra la API de WhatsApp con un modelo de lenguaje para responder consultas de clientes en lenguaje natural (Autor, año). & Se implementará en la organización y se aplicará como tratamiento al grupo experimental durante el periodo de post-prueba. & --- & --- & --- & --- & --- \\
\midrule
\textbf{VD:} Gestión de atención al cliente & EJEMPLO: conjunto de actividades con las que una organización recibe, atiende y resuelve las consultas de sus clientes (Autor, año). & Se medirá mediante el tiempo de primera respuesta y la tasa de resolución, registrados en fichas de registro antes y después de implementar el sistema. & Oportunidad de la atención & Tiempo promedio de primera respuesta (TPR), en minutos & $\displaystyle TPR = \frac{\sum_{i=1}^{n} t_i}{n}$ & De razón & Ficha de registro \\
 &  &  & Eficacia de la atención & Tasa de consultas resueltas sin intervención humana (TRA), en porcentaje & $\displaystyle TRA = \frac{CR_{auto}}{CT} \times 100$ & De razón & Ficha de registro \\
\end{longtable}}
\vspace{-10pt}
\nota{Elaboración propia.}
\end{landscape}


\section{Hipótesis de la investigación}
% GENERADO por generar.py desde tesis.yaml. No editar a mano.
\subsection{Hipótesis general}
Un sistema web con chatbot de WhatsApp e inteligencia artificial mejora significativamente la gestión de atención al cliente en una empresa de servicios.

\subsection{Hipótesis específicas}
\begin{enumerate}[label=\textbf{HE\arabic*:},leftmargin=1.27cm,itemsep=4pt]
\item Un sistema web con chatbot de WhatsApp e inteligencia artificial reduce significativamente el tiempo promedio de primera respuesta al cliente.
\item Un sistema web con chatbot de WhatsApp e inteligencia artificial incrementa significativamente la tasa de consultas resueltas sin intervención humana.
\end{enumerate}


\chapter{METODOLOGÍA}
% METODOLOGÍA (capítulo IV con hipótesis, III sin hipótesis). Lo escribe tesis-metodologia.
% Las secciones cambian entre plan y tesis; \ifmodoplan elige cuál se imprime.

\section{Diseño de investigación}
% \subsection{Tipo}, \subsection{Nivel}, \subsection{Diseño} (esquema GE: O1 X O2),
% \subsection{Enfoque}, \subsection{Método}; cada uno definido con cita y aplicado.

\section{Descripción de la metodología}
% Cómo se ejecuta el diseño: grupo, cuándo se miden O1 y O2, indicadores e instrumentos.
% Separa el desarrollo del software de la comprobación de hipótesis.
\ifmodoplan
\subsection{Etapas del desarrollo del plan de tesis}
% Etapas de la investigación y del desarrollo del sistema (p. ej. Scrum: sprint 0 ... n).
\else
\subsection{Implementación del tema de investigación}
% Desarrollo del sistema por sprints: requerimientos, arquitectura, modelo de datos,
% integraciones (WhatsApp, OpenAI), capturas y fragmentos de código en lstlisting.
\subsection{Pruebas realizadas}
% Pruebas funcionales y cómo se tomaron las mediciones pre y post de cada indicador.
\fi

\ifconhipotesis
\section{Población y muestra}
% \subsection{Población} y \subsection{Muestra}: unidad de análisis, criterios de inclusión y
% exclusión, tamaño (fórmula si es probabilístico) y tipo de muestreo.

\ifmodoplan
\section{Técnicas e instrumentos de recolección de datos}
% Técnica -> instrumento por indicador, validez y confiabilidad previstas, plan de análisis.
\else
\section{Técnicas de recolección de datos}
% Técnica usada por indicador, con cita.

\section{Instrumentos de recolección de datos}
% Instrumento por indicador (fichas en el Anexo 2).
\subsection{Validez}
% Juicio de expertos (3): claridad, pertinencia, relevancia. Resultado en el Anexo 3.
\subsection{Confiabilidad}
% Datos del sistema: trazabilidad de registros. Cuestionarios: Alfa de Cronbach > 0,7.
\fi
\fi

\chapter{CRONOGRAMA DE ACTIVIDADES DE LA TESIS}
% CRONOGRAMA (solo plan). Párrafo breve opcional antes de la tabla generada desde tesis.yaml.

% GENERADO por generar.py desde tesis.yaml. No editar a mano.
\begin{table}[H]
\caption{Cronograma de actividades de la tesis}\label{tab:cronograma}
{\footnotesize\setlength{\tabcolsep}{1pt}
\begin{tabular}{@{}L{4.6cm}*{22}{C{0.394cm}}@{}}
\toprule
\multirow{2}{*}{\textbf{Actividad}} & \multicolumn{22}{c}{\textbf{Semana}} \\
\cmidrule(l){2-23}
 & \tiny 1 & \tiny 2 & \tiny 3 & \tiny 4 & \tiny 5 & \tiny 6 & \tiny 7 & \tiny 8 & \tiny 9 & \tiny 10 & \tiny 11 & \tiny 12 & \tiny 13 & \tiny 14 & \tiny 15 & \tiny 16 & \tiny 17 & \tiny 18 & \tiny 19 & \tiny 20 & \tiny 21 & \tiny 22 \\
\midrule
Revisión de literatura y marco teórico & \cellcolor{gray!65} & \cellcolor{gray!65} & \cellcolor{gray!65} & \cellcolor{gray!65} &  &  &  &  &  &  &  &  &  &  &  &  &  &  &  &  &  &  \\[2pt]
Pre-prueba: registro de indicadores &  &  &  &  & \cellcolor{gray!65} & \cellcolor{gray!65} & \cellcolor{gray!65} & \cellcolor{gray!65} &  &  &  &  &  &  &  &  &  &  &  &  &  &  \\[2pt]
Desarrollo del sistema (sprints) &  &  &  &  & \cellcolor{gray!65} & \cellcolor{gray!65} & \cellcolor{gray!65} & \cellcolor{gray!65} & \cellcolor{gray!65} & \cellcolor{gray!65} & \cellcolor{gray!65} & \cellcolor{gray!65} & \cellcolor{gray!65} & \cellcolor{gray!65} &  &  &  &  &  &  &  &  \\[2pt]
Post-prueba: registro de indicadores &  &  &  &  &  &  &  &  &  &  &  &  &  &  & \cellcolor{gray!65} & \cellcolor{gray!65} & \cellcolor{gray!65} & \cellcolor{gray!65} &  &  &  &  \\[2pt]
Análisis estadístico y redacción de la tesis &  &  &  &  &  &  &  &  &  &  &  &  &  &  &  &  &  & \cellcolor{gray!65} & \cellcolor{gray!65} & \cellcolor{gray!65} & \cellcolor{gray!65} & \cellcolor{gray!65} \\[2pt]
\bottomrule
\end{tabular}}
\nota{Elaboración propia.}
\end{table}

\chapter{PRESUPUESTO DE LA TESIS}
% PRESUPUESTO (solo plan). Párrafo breve opcional (fuente de financiamiento) antes de la
% tabla generada desde tesis.yaml. Nunca supongas montos: cotiza.

% GENERADO por generar.py desde tesis.yaml. No editar a mano.
\begin{longtable}{@{}L{2.4cm}L{6.4cm}r r r@{}}
\caption{Presupuesto de la tesis}\label{tab:presupuesto}\\
\toprule
\textbf{Rubro} & \textbf{Descripción} & \textbf{Cant.} & \textbf{C. unit. (S/)} & \textbf{Subtotal (S/)} \\
\midrule
\endfirsthead
\toprule
\textbf{Rubro} & \textbf{Descripción} & \textbf{Cant.} & \textbf{C. unit. (S/)} & \textbf{Subtotal (S/)} \\
\midrule
\endhead
\bottomrule
\endlastfoot
Bienes & EJEMPLO: materiales de impresión y empastado & 1 & 250.00 & 250.00 \\
\multicolumn{4}{@{}r}{\textit{Subtotal bienes}} & \textit{250.00} \\
\midrule
Servicios & EJEMPLO: API de OpenAI (consumo estimado por 6 meses) & 6 & 60.00 & 360.00 \\
 & EJEMPLO: servidor en la nube (VPS) & 6 & 40.00 & 240.00 \\
\multicolumn{4}{@{}r}{\textit{Subtotal servicios}} & \textit{600.00} \\
\midrule
Otros & EJEMPLO: movilidad & 10 & 10.00 & 100.00 \\
\multicolumn{4}{@{}r}{\textit{Subtotal otros}} & \textit{100.00} \\
\midrule
\multicolumn{4}{@{}r}{Subtotal} & 950.00 \\
\multicolumn{4}{@{}r}{Imprevistos (10\,\%)} & 95.00 \\
\multicolumn{4}{@{}r}{\textbf{Total}} & \textbf{1,045.00} \\
\end{longtable}
\vspace{-10pt}
\nota{Elaboración propia. Los montos provienen de cotizaciones.}

\chapter{REFERENCIAS BIBLIOGRÁFICAS}
\printbibliography[heading=none]
\capitulosinnumero{ANEXOS}
% GENERADO por generar.py desde tesis.yaml. No editar a mano.
\begin{landscape}
\anexo{Anexo 1. Matriz de consistencia}
\noindent{\small\setstretch{1}\textbf{Título:} SISTEMA WEB CON CHATBOT DE WHATSAPP E INTELIGENCIA ARTIFICIAL PARA LA GESTIÓN DE ATENCIÓN AL CLIENTE EN UNA EMPRESA DE SERVICIOS\newline \textbf{Autor:} Apellido Apellido, Nombre Nombre\par}
{\small\setlength{\tabcolsep}{4pt}
\begin{longtable}{L{3.9cm}L{3.9cm}L{3.9cm}L{4.4cm}L{5.4cm}}
\caption{Matriz de consistencia}\label{tab:matriz-consistencia}\\
\toprule
\textbf{Problemas} & \textbf{Objetivos} & \textbf{Hipótesis} & \textbf{Variables e indicadores} & \textbf{Metodología} \\
\midrule
\endfirsthead
\multicolumn{5}{@{}l}{\itshape (continuación)}\\
\toprule
\textbf{Problemas} & \textbf{Objetivos} & \textbf{Hipótesis} & \textbf{Variables e indicadores} & \textbf{Metodología} \\
\midrule
\endhead
\bottomrule
\endlastfoot
\textbf{Problema general:}\newline ¿En qué medida un sistema web con chatbot de WhatsApp e inteligencia artificial mejora la gestión de atención al cliente en una empresa de servicios? & \textbf{Objetivo general:}\newline Determinar en qué medida un sistema web con chatbot de WhatsApp e inteligencia artificial mejora la gestión de atención al cliente en una empresa de servicios. & \textbf{Hipótesis general:}\newline Un sistema web con chatbot de WhatsApp e inteligencia artificial mejora significativamente la gestión de atención al cliente en una empresa de servicios. & \textbf{VI:} Sistema web con chatbot de WhatsApp e inteligencia artificial\newline \textbf{VD:} Gestión de atención al cliente\newline \textbf{Dimensiones:} Oportunidad de la atención; Eficacia de la atención & \textbf{Enfoque:} Cuantitativo\newline \textbf{Tipo:} Aplicada\newline \textbf{Nivel:} Explicativo\newline \textbf{Diseño:} Pre-experimental con pre-prueba y post-prueba\newline \textbf{Esquema:} GE: O1 X O2\newline \textbf{Método:} Deductivo\newline \textbf{Población:} EJEMPLO: consultas recibidas por WhatsApp en un mes típico (N = 1200).\newline \textbf{Muestra:} EJEMPLO: 30 días de registro en pre-prueba y 30 días en post-prueba.\newline \textbf{Muestreo:} No probabilístico por conveniencia\newline \textbf{Técnicas:} Observación\newline \textbf{Instrumentos:} Ficha de registro \\
\midrule
\textbf{PE1:} ¿En qué medida un sistema web con chatbot de WhatsApp e inteligencia artificial reduce el tiempo promedio de primera respuesta al cliente? & \textbf{OE1:} Determinar en qué medida un sistema web con chatbot de WhatsApp e inteligencia artificial reduce el tiempo promedio de primera respuesta al cliente. & \textbf{HE1:} Un sistema web con chatbot de WhatsApp e inteligencia artificial reduce significativamente el tiempo promedio de primera respuesta al cliente. & \textbf{Indicador:} Tiempo promedio de primera respuesta (TPR) &  \\
\midrule
\textbf{PE2:} ¿En qué medida un sistema web con chatbot de WhatsApp e inteligencia artificial incrementa la tasa de consultas resueltas sin intervención humana? & \textbf{OE2:} Determinar en qué medida un sistema web con chatbot de WhatsApp e inteligencia artificial incrementa la tasa de consultas resueltas sin intervención humana. & \textbf{HE2:} Un sistema web con chatbot de WhatsApp e inteligencia artificial incrementa significativamente la tasa de consultas resueltas sin intervención humana. & \textbf{Indicador:} Tasa de consultas resueltas sin intervención humana (TRA) &  \\
\end{longtable}}
\vspace{-10pt}
\nota{Elaboración propia.}
\end{landscape}

\clearpage
\anexo{Anexo 2. Instrumentos de recolección de datos}
% GENERADO por generar.py desde tesis.yaml. No editar a mano.
% Una ficha de registro por indicador de la variable dependiente.
\begin{table}[H]
\caption{Ficha de registro del indicador Tiempo promedio de primera respuesta}\label{tab:ficha-1}
{\small\begin{tabularx}{\textwidth}{@{}lY@{}}
\toprule
\textbf{Dimensión} & Oportunidad de la atención \\
\textbf{Indicador} & Tiempo promedio de primera respuesta (TPR) \\
\textbf{Fórmula} & $\displaystyle TPR = \frac{\sum_{i=1}^{n} t_i}{n}$ \\
\textbf{Unidad de medida} & minutos \\
\textbf{Técnica / instrumento} & Observación / Ficha de registro \\
\textbf{Medición} & Pre-prueba ( )\quad Post-prueba ( ) \\
\bottomrule
\end{tabularx}}
\medskip
{\small\begin{tabularx}{\textwidth}{@{}C{1cm}C{2.6cm}C{3.4cm}Y@{}}
\toprule
\textbf{N.\textsuperscript{o}} & \textbf{Fecha} & \textbf{TPR (minutos)} & \textbf{Observaciones} \\
\midrule
1 & & & \\[3pt]
2 & & & \\[3pt]
3 & & & \\[3pt]
4 & & & \\[3pt]
5 & & & \\[3pt]
6 & & & \\[3pt]
7 & & & \\[3pt]
8 & & & \\[3pt]
9 & & & \\[3pt]
10 & & & \\[3pt]
\bottomrule
\end{tabularx}}
\nota{Elaboración propia.}
\end{table}
\begin{table}[H]
\caption{Ficha de registro del indicador Tasa de consultas resueltas sin intervención humana}\label{tab:ficha-2}
{\small\begin{tabularx}{\textwidth}{@{}lY@{}}
\toprule
\textbf{Dimensión} & Eficacia de la atención \\
\textbf{Indicador} & Tasa de consultas resueltas sin intervención humana (TRA) \\
\textbf{Fórmula} & $\displaystyle TRA = \frac{CR_{auto}}{CT} \times 100$ \\
\textbf{Unidad de medida} & porcentaje \\
\textbf{Técnica / instrumento} & Observación / Ficha de registro \\
\textbf{Medición} & Pre-prueba ( )\quad Post-prueba ( ) \\
\bottomrule
\end{tabularx}}
\medskip
{\small\begin{tabularx}{\textwidth}{@{}C{1cm}C{2.6cm}C{3.4cm}Y@{}}
\toprule
\textbf{N.\textsuperscript{o}} & \textbf{Fecha} & \textbf{TRA (porcentaje)} & \textbf{Observaciones} \\
\midrule
1 & & & \\[3pt]
2 & & & \\[3pt]
3 & & & \\[3pt]
4 & & & \\[3pt]
5 & & & \\[3pt]
6 & & & \\[3pt]
7 & & & \\[3pt]
8 & & & \\[3pt]
9 & & & \\[3pt]
10 & & & \\[3pt]
\bottomrule
\end{tabularx}}
\nota{Elaboración propia.}
\end{table}

% Instrumentos adicionales del Anexo 2 (p. ej. un cuestionario). Las fichas de registro por
% indicador se generan desde tesis.yaml y aparecen antes de este contenido.


\end{document}
